% Junior Level 3, Week 6 -- workbook chapter (read aloud by a grown-up)
\chapterweek{6}{Hello, Python!}{Python uses words instead of blocks, and it is picky about spelling.}

\section{What you will learn}
\begin{ican}
  \item I can run a Python program in Thonny.
  \item I can change a number and type one new line.
  \item I can find a typo with the help of the red text.
\end{ican}

\section{Meet Thonny}
Thonny is the app where we write Python. The turtle is our sprite. It draws wherever it walks.

\begin{center}
%% alt: The Thonny window in simple mode. A row of buttons at the top: Load, Save, a green Run button and a red Stop button. Under it the big Editor shows a short turtle program. At the bottom the Shell is where Python talks back. On the right, a separate turtle window shows a turtle at the end of a line.
\begin{tikzpicture}[font=\small,
  lab/.style={draw=kcgray,fill=white,rounded corners=2pt,inner sep=3pt,align=center}]
  \draw[thick,rounded corners=3pt,fill=kcgray!10] (0,0) rectangle (9,6.4);
  \node[anchor=west,font=\small\bfseries] at (0.2,6.15) {Thonny};
  % toolbar with labeled buttons
  \foreach \x/\name in {0.8/Load,2.0/Save} {
    \draw[fill=white,rounded corners=1pt] (\x-0.25,5.45) rectangle (\x+0.25,5.85);
    \node[font=\scriptsize] at (\x,5.25) {\name};
  }
  \fill[kcgreen] (3.3,5.65) circle (0.24);
  \fill[white] (3.21,5.52) -- (3.21,5.78) -- (3.42,5.65) -- cycle;
  \node[font=\scriptsize\bfseries] at (3.3,5.25) {Run};
  \fill[kcvermillion] (4.5,5.65) circle (0.24);
  \fill[white] (4.41,5.56) rectangle (4.59,5.74);
  \node[font=\scriptsize] at (4.5,5.25) {Stop};
  % editor
  \draw[fill=white] (0.2,2.2) rectangle (8.8,5.0);
  \node[anchor=north west,font=\ttfamily\footnotesize\setstretch{1},align=left] at (0.35,4.9)
    {import turtle\\t = turtle.Turtle()\\t.forward(150)\\t.left(90)\\turtle.done()};
  \node[lab,fill=kcsky!20] at (6.8,3.6) {\textbf{Editor}\\my program};
  % shell
  \draw[fill=kcblack!4] (0.2,0.2) rectangle (8.8,2.0);
  \node[anchor=north west,font=\ttfamily\footnotesize] at (0.35,1.9) {Shell};
  \node[lab,fill=kcyellow!30] at (6.8,1.0) {\textbf{Shell}\\Python talks back};
  % turtle window
  \draw[thick,rounded corners=3pt,fill=white] (9.6,1.2) rectangle (13.4,4.6);
  \node[anchor=west,font=\scriptsize\bfseries] at (9.7,4.35) {Turtle window};
  \draw[very thick] (10.2,2.4) -- (12.4,2.4);
  \fill[kcgreen!70!black] (12.4,2.4) ellipse (0.18 and 0.12);
  % labels
  \node[lab] (run) at (6.6,7.1) {\textbf{Run} is the green flag};
  \draw[-{Stealth},thick] (run.west) -- (3.6,5.85);
\end{tikzpicture}
\end{center}

Every Scratch block has a word twin.
\begin{center}
\renewcommand{\arraystretch}{1.6}
\begin{tabular}{@{}l c l@{}}
\begin{scratch}[scale=0.8]\blockmove{move \ovalnum{100} steps}\end{scratch} & \textrightarrow & \texttt{t.forward(100)}\\
\begin{scratch}[scale=0.8]\blockmove{turn \turnleft{} \ovalnum{90} degrees}\end{scratch} & \textrightarrow & \texttt{t.left(90)}\\
\end{tabular}
\end{center}

\begin{exercise}[Hello, Turtle!]
Folder: \texttt{week06/ex-hello-turtle}.
\begin{enumerate}[leftmargin=1.4em,itemsep=2pt]
  \item Click \menu{Load}. Open \texttt{starter.py}.
  \item Click \menu{Run}. Watch the turtle.
  \item Change \texttt{50} to \texttt{150}.
  \item Click \menu{Save}, then \menu{Run}. The line is longer!
\end{enumerate}
\end{exercise}

\begin{exercise}[Finish the Square]
Folder: \texttt{week06/ex-finish-square}. The turtle draws 3 sides. Type side 4 under \texttt{CHANGE ME}:
\begin{pycode}
t.forward(100)
\end{pycode}
Small letters. A dot. Both brackets. Save and Run. Do you see a square?
\end{exercise}

\begin{exercise}[Typo Detective]
Python stops at a typo. It shows red text in the Shell. The \textbf{line number} says where to look. The \textbf{last line} says what is wrong. Circle the typo in each message.
\begin{center}
%% alt: Two red messages from the Thonny Shell. Bug 1: File starter.py, line 9, the line t.forwad(100), and the last line AttributeError: 'Turtle' object has no attribute 'forwad'. Bug 2: File starter.py, line 16, the line t.forward(100 with a little arrow under the bracket, and the last line SyntaxError: '(' was never closed.
\begin{tikzpicture}[
  err/.style={draw=kcvermillion,fill=kcvermillion!6,rounded corners=2pt,inner sep=5pt,align=left,
    font=\ttfamily\footnotesize,text=kcvermillion!85!black},
  tag/.style={font=\small\bfseries,anchor=south west,inner sep=1pt}]
  \node[err,anchor=north west] (b1) at (0,0) {Traceback (most recent call last):\\
    \ \ File "starter.py", line 9, in <module>\\
    \ \ \ \ t.forwad(100)\\
    AttributeError: 'Turtle' object has no attribute 'forwad'};
  \node[tag] at (b1.north west) {Bug 1};
  \node[err,anchor=north west] (b2) at ([yshift=-0.9cm]b1.south west) {\ \ File "starter.py", line 16\\
    \ \ \ \ t.forward(100\\
    \ \ \ \ \ \ \ \ \ \ \ \ \ \textasciicircum\\
    SyntaxError: '(' was never closed};
  \node[tag] at (b2.north west) {Bug 2};
\end{tikzpicture}
\end{center}
Now your grown-up hides one typo in your square. Can you find it?
\end{exercise}

\begin{grownup}[No turtle?]
The turtle window can hide behind Thonny. Click the Python icon in the Dock.
\end{grownup}

\begin{vocab}
  \vocabword{Thonny}{The app where we write and run Python.}
  \vocabword{typo}{A small spelling mistake, like \texttt{forwad}.}
  \vocabword{error message}{Red text from Python. It says which line has a problem.}
\end{vocab}

\begin{reflection}
What do you like more today? Circle one.

\begin{center}
\begin{tabular}{c@{\hspace{2cm}}c}
{\Huge\faPuzzlePiece} & {\Huge\faKeyboard[regular]}\\
blocks & words
\end{tabular}
\end{center}

How do you feel today? Circle one.

\begin{center}
\Huge\faSmile[regular]\qquad\faMeh[regular]\qquad\faFrown[regular]
\end{center}
\end{reflection}
